\documentclass[11pt]{article}

\usepackage[margin=1in]{geometry}
\usepackage{amsmath,amssymb}
\usepackage{booktabs}
\usepackage{listings}
\usepackage{xcolor}
\usepackage{tcolorbox}
\usepackage[hidelinks]{hyperref}
\usepackage{enumitem}

\definecolor{codegray}{gray}{0.95}
\definecolor{codegreen}{rgb}{0.0,0.45,0.0}
\definecolor{codeblue}{rgb}{0.0,0.0,0.65}
\definecolor{boxblue}{rgb}{0.90,0.95,1.00}
\definecolor{boxyellow}{rgb}{1.00,0.97,0.85}

\lstdefinestyle{py}{
  backgroundcolor=\color{codegray},
  basicstyle=\ttfamily\small,
  keywordstyle=\color{codeblue},
  commentstyle=\color{codegreen},
  stringstyle=\color{red!60!black},
  showstringspaces=false,
  breaklines=true,
  frame=single,
  language=Python,
  columns=fullflexible
}
\lstset{style=py}

\newtcolorbox{keyidea}{colback=boxblue,colframe=blue!50!black,title=Key idea}
\newtcolorbox{lesson}{colback=boxyellow,colframe=orange!70!black,title=What this teaches}

\title{Breaking the Kerckhoffs Geocache Cipher\\\large A step-by-step walkthrough for the curious undergraduate}
\author{}
\date{}

\begin{document}
\maketitle

\begin{abstract}
Three geocaches hide a custom polyalphabetic cipher.  The first publishes its
whole table, the second hides the table but keeps two conveniences, and the
third removes both conveniences and looks unbreakable.
This document walks through how each part was solved, slowly enough to follow
with basic linear algebra, a little statistics, and introductory algorithms.
The point is not just the answer.  It is to show how a handful of standard
ideas---modular arithmetic, linear systems over finite fields, frequency
statistics, bitmasks, and greedy algorithms---combine into a real
cryptanalysis.
\end{abstract}

\tableofcontents

\section{The setting: Kerckhoffs's principle}

Auguste Kerckhoffs argued that a cryptosystem should stay secure even if
everything about its design is public except the key.  The cache series is a
playful demonstration.  Each part uses the same basic machine:

\begin{enumerate}[label=\arabic*.]
  \item There is a table mapping the 26 plaintext letters to 26 digit strings
        called \emph{codewords}.
  \item The table is used to replace each letter by its codeword.
  \item After each codeword is used, the whole table is rotated by a shift
        value.  The rotation is cumulative, so the same letter can be encoded
        by different codewords at different positions.
\end{enumerate}

Part 1 publishes the table.  Part 2 hides the table but keeps two hints in the
ciphertext: the first digit of a codeword says how long it is, and the last
digit says how much to rotate afterwards.  Part 3 removes both hints: the
first digit is meaningless and the shift is a secret value attached to each
codeword.

\begin{keyidea}
The whole series is a lesson in \emph{known-plaintext attacks}.  A cipher can
look fearsome because the key space is huge, yet collapse the moment an
attacker knows a long stretch of plaintext.  That is exactly what happens in
part 3.
\end{keyidea}

\section{Part 1: the table is public}\label{sec:part1}

\subsection{The machine}

Let the public table be a list of 26 codewords, one per letter.  In the actual
puzzle the file alternates letters and codewords:

\begin{lstlisting}
A 89411894
B 52381
C 88527391
...
Z 4502
\end{lstlisting}

Each codeword begins with its own length: \texttt{89411894} has length 8 and
indeed starts with 8; \texttt{52381} has length 5 and starts with 5.  This is
a \emph{prefix-free} code: no codeword is a prefix of another, because the
first digit tells the decoder where the codeword ends.

After encoding a letter, the table is rotated so that every letter now maps to
the codeword that used to belong to the previous letter.  Concretely, if the
table is a list of 26 codewords, a rotation by $k$ moves each codeword down
$k$ slots (wrapping around).  The rotation amount is the \emph{last digit} of
the codeword that was just used.

\subsection{Decryption by hand}

Because the first digit gives the length, decryption is a straightforward
greedy scan:

\begin{enumerate}[label=\arabic*.]
  \item Read the first digit; that is the codeword length.
  \item Slice out the codeword.
  \item Look up its plaintext letter in the current table.
  \item Rotate the table by the codeword's last digit.
  \item Repeat from the new position.
\end{enumerate}

The decrypted message contains the final coordinates:

\begin{center}
\texttt{THE CACHE IS LOCATED AT NORTH THIRTY NINE DEGREES TEN POINT FOUR EIGHT TWO \dots}
\end{center}

\begin{lesson}
Greedy decoding works because the code is prefix-free.  The first digit is an
explicit length field, so the parser never has to backtrack.  This is the same
idea as length-prefixed strings in binary formats.
\end{lesson}

\section{Part 2: the table is hidden}\label{sec:part2}

\subsection{What is still visible}

Two things survive:

\begin{itemize}
  \item The first digit of each codeword still says how long it is, so the
        ciphertext can be split into codewords by a simple scan.
  \item The last digit still says how far the table rotates after that
        codeword, so the cumulative shift before every position is known.
\end{itemize}

Let $c_i$ be the codeword at position $i$, let $\mathrm{pos}(c_i)$ be its
position in the \emph{original} table, and let $S_i$ be the cumulative
rotation before position $i$.  Then

\begin{equation}\label{eq:p2}
\boxed{\;p_i \;=\; \bigl(\mathrm{pos}(c_i) + S_i\bigr)\bmod 26\;}
\end{equation}

where $p_i$ is the $i$-th plaintext letter as a number $0,\dots,25$.  The
shifts $S_i$ are known from the ciphertext.  The only unknown is the
permutation $\mathrm{pos}$.

\subsection{The search space and a smarter move}

There are $26!$ possible tables---about $4\times 10^{26}$.  Brute force is
hopeless.  But the problem has a very useful shape: if we guess a table, we can
decrypt the whole message and ask, ``does this look like English?''

That is a classic optimisation setup:

\begin{itemize}
  \item \textbf{State:} a permutation of the 26 table positions.
  \item \textbf{Score:} how English-like the decrypted text is.
  \item \textbf{Move:} swap two entries of the permutation.
\end{itemize}

\subsection{Scoring English}

A \emph{quadgram} is a sequence of four letters, e.g.\ \texttt{THEM}.  The
repository includes counts of quadgrams from a large English corpus.  If a
candidate decryption contains quadgram $q$, we add $\log P(q)$ to its score.
For a text of length $n$ the total score is

\begin{equation}
\mathrm{score}(p_1\dots p_n) \;=\; \sum_{i=1}^{n-3} \log P(p_i p_{i+1} p_{i+2} p_{i+3}).
\end{equation}

Logs are used for two reasons: they turn a product of probabilities (which
would underflow) into a sum, and the sum is cheap to update when only a few
letters change.

\begin{lesson}
Language models are just probability distributions over short letter
sequences.  Turning a product into a sum of logs is a standard numerical trick,
and it makes incremental rescoring possible.
\end{lesson}

\subsection{Simulated annealing}

Simulated annealing is a randomised hill-climber.  At each step we propose a
swap.  If the score improves we accept it.  If it gets worse we accept it with
probability

\begin{equation}
\exp\!\left(\frac{\Delta}{T}\right),
\end{equation}

where $\Delta$ is the score change and $T$ is a \emph{temperature} that starts
high (explore) and slowly cools (exploit).  This lets the search escape local
optima early on and settle into a good solution later.

The implementation is careful about speed.  When we swap the entries for
symbols $a$ and $b$, only positions whose codeword is $a$ or $b$ change their
plaintext letter, and only quadgrams overlapping those positions change.  A
precomputed list of positions per symbol (a hash map from symbol to positions)
makes each move cost roughly the number of affected positions rather than the
length of the text.

\begin{lesson}
The algorithmic idea is \emph{incremental evaluation}: never recompute more
than the move actually changed.  This is the same reason balanced trees and
hash maps matter in algorithmic problem solving.
\end{lesson}

The annealing recovers the table and decrypts the message.  It also reveals the
coordinate instruction:

\begin{center}
\texttt{TO SOLVE SWAP AND DOUBLE THE LAST PARTS OF THE COORDS}
\end{center}

Applying it to the posted coordinates $N\,39^\circ\,10.227$, $W\,076^\circ\,
51.187$ means swapping the decimal minutes ($227/187 \to 187/227$) and
doubling them ($187/227 \to 374/454$), giving the confirmed final coordinate
$N\,39^\circ\,10.374$, $W\,076^\circ\,51.454$.  Applying the same recipe to
part 1's numbers is a trap --- the checker rejects it.

\section{Part 3: both conveniences removed}\label{sec:part3}

Part 3 changes two rules:

\begin{itemize}
  \item the first digit no longer gives the codeword length;
  \item the last digit no longer gives the shift: every codeword has a secret
        shift value in $0,\dots,25$.
\end{itemize}

Without the length prefix, the ciphertext cannot even be split into codewords
without knowing the dictionary.  Without the shift digits, the cumulative
shift is part of the key.  This section is the heart of the document.

\subsection{Recovering the codeword dictionary}\label{sec:dict}

The plaintext is a concatenation of codewords, but we do not know where the
boundaries are.  However, the code is still prefix-free, and the codewords
repeat often: a codeword encodes a fixed plaintext letter (up to the table
rotation), so a common letter's codeword appears many times.

\paragraph{A statistical argument.} Suppose the dictionary were replaced by
random digit strings and the text were random.  The expected number of
occurrences of a fixed length-$L$ string in $N$ digits is about $N/10^L$.  For
$N=3377$ and $L=9$ that is $3.4\times10^{-6}$---you would essentially never see
a 9-digit string repeat.  Yet the recovered length-9 codewords appear 18--26
times each.  That is overwhelming evidence that they are real codewords rather
than coincidences.

\paragraph{The recovery procedure.}
\begin{enumerate}[label=\arabic*.]
  \item Count occurrences of every substring of length $2,\dots,9$ with a hash
        map.  This is $O(9N)$ work.
  \item Keep the most frequent maximal strings; these are codeword candidates.
  \item Greedily segment the whole ciphertext with the candidates, longest
        first.  A prefix-free set makes the segmentation unambiguous.
  \item Verify that the set tiles the entire ciphertext with no leftovers and
        that no candidate is a prefix of another.
\end{enumerate}

The result is 26 codewords, e.g.

\begin{lstlisting}
243255912  222146337  946294413  80920926  63388773
37638005   6103003    0273351    2167948   8621836
...
860        781        418        0648      30    93
\end{lstlisting}

\subsection{The key and the model}

Now the decryption rule is

\begin{align}
p_i &= \bigl(\tau[c_i] + S_i\bigr)\bmod 26,\label{eq:p3}\\
S_{i+1} &= \bigl(S_i + \omega[c_i]\bigr)\bmod 26,\label{eq:shift}
\end{align}

where

\begin{itemize}
  \item $c_i$ is the (now known) codeword index at position $i$;
  \item $\tau$ is a permutation of $0,\dots,25$ (the hidden table order);
  \item $\omega[c]$ is the secret rotation after codeword $c$.
\end{itemize}

There are 26 values of $\tau$ and 26 values of $\omega$: 52 unknowns.  The
naive key space is $26!\cdot 26^{26}$, which is astronomically large.

\subsection{The linear structure that makes it solvable}\label{sec:linear}

The secret shifts look intimidating, but the cipher is \emph{linear}.  Subtract
the rule at consecutive positions:

\begin{align}
p_{i+1}-p_i
&= \bigl(\tau[c_{i+1}]+S_{i+1}\bigr) - \bigl(\tau[c_i]+S_i\bigr)\notag\\
&= \tau[c_{i+1}] - \tau[c_i] + \omega[c_i] \pmod{26}.
\label{eq:diff}
\end{align}

Everything on the right is a sum of key values with coefficients $0,\pm1$.
If we know a stretch of plaintext---a \emph{crib}---then
equation~\eqref{eq:diff} gives one linear equation per adjacent pair of
letters, plus one absolute equation for the first letter.  That is a linear
system in the 52 unknowns over the ring $\mathbb{Z}/26\mathbb{Z}$.

\begin{keyidea}
A crib of roughly 100 letters gives more than 52 equations.  The system is
overdetermined and usually has a unique solution for the key, even though the
key space looked like $26!\cdot 26^{26}$.  The apparent hardness was the
missing plaintext, not the key size.
\end{keyidea}

\subsection{Solving the linear system}

The modulus factors as $26 = 2\cdot 13$, and both 2 and 13 are prime.  By the
Chinese remainder theorem, solving modulo 26 is the same as solving modulo 2
and modulo 13 separately and combining the answers.  Over a prime field,
ordinary Gaussian elimination works: we can divide by any nonzero coefficient.

\begin{enumerate}[label=\arabic*.]
  \item Build the equations from the crib.
  \item Perform Gaussian elimination modulo 2 and modulo 13.
  \item If either system is inconsistent, the crib is wrong (or misplaced).
  \item Otherwise read off free variables, enumerate their small number of
        possibilities, and keep only keys whose $\tau$ is a permutation.
  \item Decrypt and score with the quadgram model.
\end{enumerate}

\begin{lesson}
Finite fields, Gaussian elimination, and the Chinese remainder theorem are not
just abstract algebra.  They are exactly the tools that turn a huge discrete
search into a small linear-algebra computation.
\end{lesson}

\subsection{Finding the crib: hints and a two-block search}

Part 3 has no public table and no checker, so the crib is the whole game.  The
hint block eventually supplied:

\begin{tcolorbox}[colback=boxyellow,colframe=orange!70!black]
[Puzzle 4] The two phrases \texttt{North Thirty Nine Degrees Ten Point} and
\texttt{West Seventy Six Degrees Fifty One Point} both appear somewhere in the
plaintext.  No other assumptions about the plaintext need be made.
\end{tcolorbox}

The earlier hint (``only the seconds differ'') only guarantees that the final
coordinate shares its degrees and minutes with the posted one.  Hint 4 is what
gives \emph{exact} plaintext.

Write the two phrases as

\begin{align*}
A &= \texttt{NORTHTHIRTYNINEDEGREESTENPOINT} &&(30\text{ letters}),\\
B &= \texttt{WESTSEVENTYSIXDEGREESFIFTYONEPOINT} &&(34\text{ letters}).
\end{align*}

Together $A$ and $B$ are 64 known letters, more than the 52 unknowns.  The
catch is that we do not know their offsets, and they are \emph{not} adjacent
and not in the same sentence position as in part 1.

\paragraph{Searching all offset pairs.} There are about $550$ offsets for $A$
and $545$ for $B$, so roughly $3\times10^5$ pairs.  For each pair we build the
combined linear system over $\mathbb{Z}/2\mathbb{Z}$ and test consistency.

Why check modulo 2 first?  Because each equation can be packed into a Python
integer: 52 bits for the coefficients and one bit for the right-hand side.  A
53-bit integer fits in a machine word, and elimination on such rows is fast.
The entire $\sim3\times10^5$-pair search runs in about one second in ordinary
Python.

\begin{lstlisting}
MASK = (1 << 52) - 1          # 52 coefficient bits; bit 52 is the RHS

def consistent(a_pivots, b_rows):
    local = [0] * 53
    for row in b_rows:
        r = row
        while True:
            coeff = r & MASK
            if coeff == 0:
                if (r >> 52) & 1:
                    return False          # reduced to 0 = 1
                break
            c = coeff.bit_length() - 1    # highest set coefficient
            p = a_pivots[c] or local[c]
            if p:
                r ^= p                    # GF(2) addition is XOR
            else:
                local[c] = r
                break
    return True
\end{lstlisting}

\begin{lesson}
Choosing the right representation is an algorithmic superpower.  ``Vectors
over $\mathrm{GF}(2)$'' become bitmasks, XOR becomes vector addition, and
Gaussian elimination becomes a tight loop over machine words.
\end{lesson}

The search finds one pair that survives modulo 2 \emph{and} modulo 13: $A$ at
symbol 46, $B$ at symbol 114.  The system still has 6 free dimensions modulo 2
and 6 modulo 13, so the key is not yet unique.

\paragraph{Reading the gap.}  Decrypting with any convenient representative of
the solution space (free variables set to zero) exposes the text between the
two phrases.  It begins

\begin{center}
\texttt{FIVEHUNDREDTHIRTEENINTHEOTHERDIRECTION}
\end{center}

That is the latitude seconds \texttt{FIVEHUNDREDTHIRTEEN} = 513 followed by
the connector \texttt{INTHEOTHERDIRECTION}.  Adding this 38-letter gap as a
final crib pins the key uniquely.  The recovered key is

\begin{align*}
\tau &= [9,16,15,11,2,12,23,19,10,18,5,1,0,17,25,3,24,7,13,4,22,20,8,6,21,14],\\
\omega &= [17,13,5,21,21,17,2,5,17,10,4,14,0,1,11,18,14,24,1,5,1,24,11,5,18,19].
\end{align*}

\subsection{The plaintext and the final answer}

The decrypted message reads (spaced out for readability):

\begin{quote}
WOW IM IMPRESSED. I THINK THIS PUZZLE IS PRETTY HARD. GO TO NORTH THIRTY NINE
DEGREES TEN POINT FIVE HUNDRED THIRTEEN IN THE OTHER DIRECTION WEST SEVENTY
SIX DEGREES FIFTY ONE POINT SIX EIGHT TWO. IM REALLY CURIOUS IF YOU CAN FIND
A SOLUTION TO THIS GENERALIZATION THAT DOESNT RELY ON CRIBS. THE ONLY SOLUTION
I FOUND USED THE WORD THE AS THE MOST COMMON TRIGRAM STARTING WITH EACH
CODEWORD. EVEN THEN IT TOOK MILLIONS OF LETTERS, WASNT CONSISTENT, AND STILL
NEEDED A LARGE EXHAUST. I REALLY THOUGHT I COULD USE THE FACT THAT THERE IS A
STATISTICAL DIFFERENCE IN HOW FAR AWAY THE SECOND LETTER IS CIRCULARLY FROM
THE FIRST IN ALL BIGRAMS REGARDLESS OF FIRST LETTER, BUT IT DIDNT WORK. CACHE
IS MAYBE KNEE HIGH IN TREE.
\end{quote}

The coordinate sentence gives

\begin{equation*}
\boxed{\;N\,39^\circ\,10.513' \quad W\,076^\circ\,51.682'\;}
\end{equation*}

and the last sentence tells the finder where to look: knee high in a tree.
The posted and final coordinates differ only in the seconds, exactly as the
``only the seconds differ'' hint promised.

\section{What was tried and why it failed}\label{sec:failed}

It is worth recording the dead ends, because they show what makes the puzzle
genuinely hard and how the successful method differs.

\subsection{Coordinate wording brute force}

The earliest attempts assumed the part 3 plaintext reused part 1's coordinate
sentence verbatim, with only the numbers changed.  We generated many wordings
of ``THE CACHE IS LOCATED AT NORTH \dots'' with plausible connectors
(``YOULL ALSO NEED'', ``AND WEST'', etc.), number styles (digit-by-digit vs.\
number words), and coordinate candidates, and slid each one over every offset.

\emph{Why it failed.}  The hypothesis was too specific.  Part 3 replaces the
``THE CACHE IS LOCATED AT'' prefix, uses ``IN THE OTHER DIRECTION'' instead of
``YOULL ALSO NEED'', and the two coordinate phrases are not adjacent.  Every
one of those differences makes the crib text wrong, and a wrong crib is
rejected by the linear system.  No amount of sliding fixes a wrong \emph{word}.

\subsection{Frequency and difference statistics without a crib}

Several crib-free attacks were tried:

\begin{itemize}
  \item \textbf{Letter-frequency annealing.} Optimise the key so that the
        decrypted letter frequencies match English.
  \item \textbf{Bigram-frequency annealing.} Same idea with two-letter
        frequencies.
  \item \textbf{Letter-difference distribution.} In any English bigram
        $xy$, the circular difference $y-x \bmod 26$ has its own distribution.
        Group the $676$ codeword pairs by this difference and optimise the key
        to match the expected distribution.
  \item \textbf{Statistical recovery of the pair-difference table.} Because
        $p_{i+1}-p_i$ depends only on the codeword pair $(c_i,c_{i+1})$, the
        $26\times26$ table of pair differences is a Latin square.  One can try
        to recover it from the observed pair counts using an assignment
        (Hungarian) algorithm.  With only about $580$ pairs and counts of
        $1$--$5$ per cell, however, the counts are far too noisy to pin the
        table down.
  \item \textbf{Letter-difference trigrams.} Build the distribution of three
        consecutive differences from a quadgram corpus and anneal against it.
\end{itemize}

\emph{Why they failed.}  All of these objectives throw away the order of the
text.  They are low-dimensional summaries, so many wrong keys produce
statistics that look plausible.  The resulting optimisation landscape is full
of local optima, and randomised search gets stuck.  On synthetic data with a
known key, none of these methods recovered the key.  The author's own note in
the plaintext says the same thing happened to him: even using \texttt{THE} as a
crib took millions of letters, was inconsistent, and still needed a large
exhaustive search.  The statistical signal is real but far too weak on a
579-letter message.

\subsection{Short cribs and underdetermined systems}

Early attempts also tested short phrases from the hints (for example
\texttt{ONLYTHESECONDSDIFFER}) at every offset.  Each phrase is shorter than
the 52 unknowns, so the linear system has more unknowns than equations: it is
\emph{underdetermined}, and almost every placement is consistent.  The
consistency test therefore filters nothing, and the solver can only report
that the solution space is too large.

The lesson is a dimension count.  For a unique key you need at least as many
independent equations as unknowns, and in practice a comfortable margin.  A
30-letter phrase gives 30 equations against 52 unknowns; two 30--40 letter
phrases together give over 60 and become useful.

\subsection{Why the crib route works so much better}

The crib does not summarise the text; it gives actual letters.  Each adjacent
pair is a hard linear constraint.  With enough constraints the key is forced,
and linear algebra finds it exactly.  Cribs turn a statistical search into a
deterministic one.

\begin{lesson}
When a system is linear, one long known plaintext beats any amount of
statistical cleverness.  Known-plaintext attacks are powerful precisely
because they use exact information instead of aggregate information.
\end{lesson}

\section{Putting it all together}

The table below maps each stage of the solve to the standard topics it uses.

\begin{center}
\renewcommand{\arraystretch}{1.3}
\begin{tabular}{@{}p{0.30\textwidth}p{0.60\textwidth}@{}}
\toprule
\textbf{Step} & \textbf{Concept} \\
\midrule
Part 1 decryption & Prefix-free codes, greedy parsing, modular rotation \\
Part 2 table recovery & Permutations, simulated annealing, $n$-gram language models, incremental scoring \\
Dictionary recovery & Hash maps, substring counting, greedy segmentation, prefix trees, birthday-style statistical reasoning \\
Key model & Modular arithmetic, linear equations over a ring \\
Solving the key & Gaussian elimination, finite fields $\mathrm{GF}(2)$ and $\mathrm{GF}(13)$, Chinese remainder theorem, free variables \\
Offset search & Bitmasks, XOR as vector addition, bit parallelism, complexity analysis \\
Gap reveal & Partial decryption, solution-space representatives \\
Final answer & Coordinate arithmetic, the actual cache hunt \\
\bottomrule
\end{tabular}
\end{center}

\begin{keyidea}
The puzzle is a compressed tour of applied mathematics and computer science.
It starts as a data-structures exercise (prefix codes and hash maps), becomes a
statistics-and-optimisation exercise (language models and annealing), then
turns into linear algebra over finite fields, and finishes with a bit-level
search optimisation.  None of the individual pieces is exotic.  The difficulty
is seeing how they fit together---and, in part 3, finding the one crib that
lets the linear structure do its work.
\end{keyidea}

\section*{Further reading}
\begin{itemize}
  \item Kerckhoffs's principle (Wikipedia) for the security philosophy.
  \item Practical Cryptography's quadgram page for the language model used
        here.
  \item Any introductory linear algebra text for Gaussian elimination; any
        abstract algebra text for finite fields and the Chinese remainder
        theorem.
  \item The repository files \texttt{src/part1.py}, \texttt{src/part2.py},
        and \texttt{src/part3.py} for the actual code, and
        \texttt{POSTMORTEM.md} for the short technical summary.
\end{itemize}

\end{document}
